\begin{titlepage}
\thispagestyle{empty}
\begin{tikzpicture}[remember picture,overlay]
\fill[Ink] (current page.north west) rectangle ([xshift=13mm]current page.south west);
\draw[Accent,line width=1.1pt] ([xshift=26mm,yshift=-35mm]current page.north west)--([xshift=-22mm,yshift=-35mm]current page.north east);
\node[anchor=south east,rotate=90,text=white,font=\sffamily\small] at ([xshift=9mm,yshift=22mm]current page.south west) {LECTURE NOTES IN MATHEMATICS \quad 1553};
\draw[Ink!12,line width=0.5pt] ([xshift=-60mm,yshift=56mm]current page.south east) circle (23mm);
\draw[Ink!12,line width=0.5pt] ([xshift=-60mm,yshift=56mm]current page.south east) ellipse (36mm and 15mm);
\draw[Ink!12,line width=0.5pt] ([xshift=-60mm,yshift=56mm]current page.south east) ellipse (15mm and 36mm);
\end{tikzpicture}
\vspace*{28mm}
{\sffamily\small\color{Accent}C.I.M.E.\quad 特伦托，1991}\par
\vspace{16mm}
{\fontsize{35}{46}\selectfont\bfseries\color{Ink}算术代数几何}\par
\vspace{8mm}
{\Large\color{Ink}Arithmetic Algebraic Geometry}\par
\vspace{10mm}
{\large 1991 年 C.I.M.E. 第二期暑期学校讲义}\par
\vspace{16mm}
{\large Jean-Louis Colliot-Thélène\\[5pt]Kazuya Kato\\[5pt]Paul Vojta}\par
\vspace{9mm}
{\color{Muted}Edoardo Ballico\quad 编}\par
\vfill
{\color{Muted}中文译本·译者补注版\quad ·\quad XeLaTeX 排印}\par
\vspace{5mm}
{\small\color{Muted}原版：Springer-Verlag，1993\quad｜\quad 数学讲义丛书第 1553 卷}\par
\vspace{15mm}
\end{titlepage}

\chapter*{原书出版信息}
\addcontentsline{toc}{chapter}{原书出版信息}
\noindent\textbf{数学讲义丛书（Lecture Notes in Mathematics）1553}

丛书编委：A. Dold（海德堡）、B. Eckmann（苏黎世）、F. Takens（格罗宁根）。

佛罗伦萨 C.I.M.E. 基金会分系列；顾问：Roberto Conti。

本书收录 1991 年 6 月 24 日至 7 月 2 日在意大利特伦托举行的国际数学暑期中心（Centro Internazionale Matematico Estivo，C.I.M.E.）第二期讲习班讲义。

\subsection*{作者与编者}
\begin{description}
\item[Jean-Louis Colliot-Thélène]巴黎南大学数学系，425 号楼，F-91005 Orsay Cedex，法国。
\item[Kazuya Kato（加藤和也）]东京工业大学数学系，Oh-Okayama，Meguro-ku，Tokyo，日本。
\item[Paul Vojta]加利福尼亚大学伯克利分校数学系，Berkeley，CA 94720，美国。
\item[Edoardo Ballico（编）]特伦托大学数学系，38050 Povo，Trento，意大利。
\end{description}

\noindent 1991 年数学主题分类：14C15，14C25，14C35，19E15，14G40，11J99，11M99，11G99，11R23，14M20。

\noindent ISBN 3-540-57110-8\quad Springer-Verlag，柏林、海德堡、纽约。\\
ISBN 0-387-57110-8\quad Springer-Verlag，纽约、柏林、海德堡。

\subsection*{原书版权声明译文}
本书受版权保护。全部或部分内容的各项权利均予保留，特别包括翻译、重印、插图再利用、朗诵、广播、缩微胶片或其他形式的复制，以及在数据库中储存等权利。只有在 1965 年 9 月 9 日《德国著作权法》现行条文允许的范围内，才可复制本书或其中部分内容；使用时始终须取得 Springer-Verlag 的许可。违反规定者将依德国著作权法受到追究。

\noindent © Springer-Verlag Berlin Heidelberg 1993。德国印刷。\\
印刷与装订：Druckhaus Beltz，Hemsbach/Bergstr.\\
46/3140-543210；使用无酸纸印刷。

原文件另附 2008 年版勘误页，其 ISBN 为 978-3-540-57110-0，版权标识为 © Springer-Verlag Berlin Heidelberg 2008。该勘误及补充文献完整译于书末。

\chapter*{译本使用说明}
\addcontentsline{toc}{chapter}{译本使用说明}
本译本按原书三篇讲义的顺序编排，保留定理、命题、引理、公式及文献的原有编号。页边“原书”标记对应原书印刷页码；译本页码独立编排。文中叙述及“最近”“尚未发表”等时间表述以原书写作年代为准。

技术名词采用“中文（原文术语）”的形式：第一篇括注法语，后两篇括注英语。单击术语可跳至文末中英法对照表；表中“正文位置”列的页码可返回相应出现位置。同页有多次出现时，重复页码按出现先后排列，分别指向各处。译者补充的第三语言对应词仅供对照。

原文件包含两页无可提取文字层的后补勘误，译本将其单列于书末，并在受实质性修正影响的正文处加注。原书参考文献的作者名、题名和出版信息保留原文，便于检索；参与者名单中的人名与通信地址也予以保留。

本次补注版依据另行收到的勘误说明，在相应正文处加入编号 1--29 的“译者补注”，并附可往返跳转的索引。补注保留该说明提出的修正式、必要假设和理由；定理 6.4 的问题明确记为证明范围的缺口。原书表述及作者正式勘误保留，供读者对照。

\chapter*{序言}
\addcontentsline{toc}{chapter}{序言}
1991 年 6 月 24 日至 7 月 2 日，C.I.M.E. 举办的“算术代数几何”讲习班在意大利特伦托的 Villa Madruzzo 举行。

讲习班安排了以下系列讲座：
\begin{itemize}
\item Jean-Louis Colliot-Thélène：挠代数闭链与代数 $K$ 理论；
\item Kazuya Kato：经由 $B_{\dR}$ 研究 Hasse--Weil $L$ 函数的岩泽理论方法；
\item Christophe Soulé：Arakelov 几何；
\item Paul Vojta：算术几何在丢番图逼近中的应用。
\end{itemize}
此外，参加者还作了以下专题报告：
\begin{itemize}
\item Dan Abramovich：阿贝尔簇及 Jacobian 簇的子簇；
\item Luca Barbieri Viale：通过上同调理论研究双有理不变量；
\item Torsten Ekedahl：a）中国剩余定理的一个无限版本；b）具有一般分歧类型的扩张的密度；
\item Frans Oort：a）阿贝尔簇的复乘提升；b）Coleman 关于复数域上具有复乘的 Jacobian 簇的一个猜想；c）Newton 多边形与阿贝尔簇；
\item Angelo Vistoli：双变相交理论与 Alexander 对偶；
\item Christoph Wirsching：$G(2,4)$ 上的 Quillen 度量。
\end{itemize}
本卷收入其中三个系列讲座的扩充稿。Soulé 讲座的内容可参见 Ch. Soulé、D. Abramovich、J. F. Burnol 与 J. K. Kramer 最近出版的专著 \emph{Arakelov Geometry}，Cambridge Studies in Advanced Mathematics 33，Cambridge University Press，1992。

四位主讲人都煞费苦心，从有志钻研而尚非本领域专家的研究生所能理解的层次讲起，逐步引向研究前沿，乃至前沿之外。在本次 C.I.M.E. 讲习班之前、期间和之后，他们还以各种方式为听众提供了进一步的帮助。仅举一例：J.-L. Colliot-Thélène 早在 1990 年便编出一份长篇参考书目，既含入门资料，也含高深文献，并附有详细的内容说明；这份书目曾广泛流传。收入本书的三篇讲义，其作者在撰写时也付出了同样的心力。

在会议筹备和本书编印过程中，许多人在数学和具体事务两方面都给予了宝贵帮助，包括所有讲授者和参加者。我谨向他们致以特别的谢意。

\begin{flushright}Edoardo Ballico\end{flushright}
